% Technical / engineering report template.
% Compile: latexmk -pdf report.tex   (two passes for TOC/refs, handled automatically)
\documentclass[11pt,a4paper]{report}

\usepackage[margin=2.5cm]{geometry}
\usepackage{amsmath,amssymb}
\usepackage{graphicx}
\usepackage{booktabs}
\usepackage{siunitx}
\usepackage{caption}
\usepackage[hidelinks]{hyperref}

\sisetup{per-mode=symbol}

\title{\textbf{Technical Report Title}\\[0.5em]
       \large Subsystem or Project Name}
\author{Author Name\\ \small Organization / Team}
\date{\today}

\begin{document}

% ---------- Title page ----------
\begin{titlepage}
  \centering
  \vspace*{3cm}
  {\Huge\bfseries Technical Report Title\par}
  \vspace{1em}
  {\Large Subsystem or Project Name\par}
  \vspace{3cm}
  {\large Author Name\par}
  {\large Organization / Team\par}
  \vfill
  {\large Document ID: TR-2026-001 \quad Rev.\ A\par}
  {\large \today\par}
\end{titlepage}

\tableofcontents

\chapter{Introduction}
\section{Purpose}
State what this report documents and for whom.

\section{Scope}
What is covered, what is explicitly out of scope.

\chapter{System Description}
\section{Overview}
Reference figures as Figure~\ref{fig:block}. Replace the placeholder
rule with \verb|\includegraphics[width=0.8\textwidth]{block.pdf}|.
\begin{figure}[htb]
  \centering
  \rule{0.6\textwidth}{3cm}% placeholder
  \caption{System block diagram (placeholder).}
  \label{fig:block}
\end{figure}

\section{Key Parameters}
Units via siunitx: supply voltage \SI{3.3}{\volt}, clock
\SI{125}{\mega\hertz}, power \SI{1.2}{\watt}, data rate
\SI{10}{\giga\bit\per\second}.

\begin{table}[htb]
  \centering
  \caption{Electrical characteristics.}
  \label{tab:specs}
  \begin{tabular}{l S[table-format=3.1] l}
    \toprule
    {Parameter} & {Value} & {Unit} \\
    \midrule
    Supply voltage    & 3.3   & \si{\volt} \\
    Active current    & 210.0 & \si{\milli\ampere} \\
    Clock frequency   & 125.0 & \si{\mega\hertz} \\
    Sleep current     & 0.8   & \si{\micro\ampere} \\
    \bottomrule
  \end{tabular}
\end{table}

\chapter{Analysis}
\section{Model}
\begin{equation}
  V_\mathrm{out} = V_\mathrm{in} \frac{R_2}{R_1 + R_2}
  \label{eq:divider}
\end{equation}
Equation~\eqref{eq:divider} gives the unloaded divider output.

\section{Results}
Summarize measurements; refer to Table~\ref{tab:specs}.

\chapter{Conclusions and Recommendations}
Findings first, then concrete recommendations as a numbered list.
\begin{enumerate}
  \item Recommendation one.
  \item Recommendation two.
\end{enumerate}

\appendix
\chapter{Test Setup Details}
Instrument list, firmware versions, environmental conditions.

\chapter{Raw Data}
Long tables and logs go here, out of the main narrative.

\end{document}
